\sectionguard
\section{Liturgical and Magisterial Reception}

\subsection{The cult before Rome, 1555--1666}

The devotion is documented before any Roman act touches it. By the
mid-1550s the hermitage at Tepeyac is a functioning shrine with alms,
pilgrimage, reported cures, and---as the 1556 quarrel shows---enough
prominence to divide an archbishop from a Franciscan provincial
(\S3.2). Testaments, maps, and viceregal correspondence attest it
across the second half of the century (\S2.3). What the sixteenth
century does not produce is a juridical act about the events.

The first attempt to obtain one is the petition of 1666. The chapter
of the collegiate church gathered the \work{Informaciones
jur\'idicas}; a Latin relation was drawn up and presented to the
Sacred Congregation of Rites, together with the postulates of the
city of Mexico, of the dean and chapter \term{sede vacante}, and of
the religious bodies, asking Alexander~VII for a proper Office and
Mass in honour of the apparition. It did not
succeed.\endnote{The purpose and addressees of the 1666 instruments
are stated in the prologue of \work{Colecci\'on de obras y opusculos}
(Madrid, 1785), read 2026-07-25 at
\url{https://archive.org/details/coleccindeobrasy00unkn}, which
describes Anastasio Nicoseli's Italian relation of 1781 as drawn
``de una Relaci\'on Latina, que se present\'o en la Sagrada
Congregaci\'on de Ritos, con los Postulados de la Nobil\'isima Ciudad
de M\'exico, del muy Ilustre y Venerable Dean y Cabildo Sede-vacante,
y de todos los Cuerpos Religiosos, \'a fin de que el Sumo Pont\'ifice
Alejandro VII concediese Oficio propio y Misa en honor de esta
Aparici\'on.'' The same work records that the 1666 instruments were
supplemented by fresh diligences under the abbot and chapter installed
on 22 October 1750. That the seventeenth-century petition failed is
inferred from the fact that the concession dates from 1754 and from
the eighteenth-century sources' own account of the renewed effort; no
act of refusal was located.}

\subsection{The eighteenth century: patronage and the proper Office}

Two acts belong here, and one of them is better attested by a much
later document than by any edition of its own text.

In 1737, in the circumstances of a devastating epidemic, the republic
took Our Lady of Guadalupe as its patroness by a sworn vow. In 1754
Benedict~XIV confirmed that devotion by his authority. Both facts are
stated by the Holy See itself, in Pius~X's letter of 16 June 1910
granting indulgences for the centenary renewal of the same oath.

\begin{operativetext}{Pius~X, apostolic letter granting indulgences to
the faithful of the Mexican Republic, 16 June 1910, extract (Latin, as
printed in \work{Acta Apostolicae Sedis})}
\ldots\ proposuit, ut a Mexico iusiurandum renovaretur, quo inde ab
anno MDCCXXXVII se suaeque Deiparae Virginis a Guadalupe Patrocinio
respublica credidit atque commisit. Cum vero huiusmodi religionem
fel.\ rec.\ Benedictus XIV Decessor Noster anno MDCCLIV auctoritate
sua confirmaverit\ldots\endnote{\work{Acta Apostolicae Sedis} 2
(1910), p.~569, read 2026-07-25 in the Holy See's OCR scan
(\url{https://www.vatican.va/archive/aas/documents/AAS-02-1910-ocr.pdf}),
page located and read at first hand. Editorially: ``\ldots proposed
that Mexico renew the oath by which, from the year 1737, the republic
entrusted and committed itself and its own to the patronage of the
Virgin Mother of God of Guadalupe. And since Benedict~XIV, Our
predecessor of happy memory, confirmed this devotion by his authority
in the year 1754\ldots'' Quoted within ordinary scholarly limits with
citation.}
\end{operativetext}

Beyond that, precision fails, and the failure should be recorded
rather than covered. The eighteenth-century Guadalupan literature
names the act as the brief \work{Non est equidem} of May 1754 and
describes the concession as a proper Office and Mass under the most
solemn rite, first class with octave, together with the apostolic
confirmation of the universal patronage. This study did not reach the
Latin text of that brief in any edition, and a bibliographic anomaly
turned up in the attempt.\endnote{The incipit and year are attested in
\work{Colecci\'on de obras y opusculos} (Madrid, 1785): ``\ldots como
se registra en el Breve \term{Non est equidem} su data en Roma en
\ldots\ de Mayo de 1754'' (the day of the month is garbled in the OCR
of the copy read and is not asserted here), with the concession
described elsewhere in the same volume as ``Rezo propio con el mas
solemne Rito de primera Clase y Octava'' and as ``la Apost\'olica
confirmaci\'on del universal Patronato''; read 2026-07-25 at
\url{https://archive.org/details/coleccindeobrasy00unkn}. Garc\'ia
Icazbalceta, \work{Carta}, \S6, independently dates ``la Bula de la
concesi\'on del patronato'' to 1754. \emph{Recorded limits and
anomaly.} (1)~No Latin text of \work{Non est equidem} concerning
Guadalupe was reached: the Mechelen \work{Bullarium} volume covering
1754 is not among the volumes digitized at the Internet Archive.
(2)~In the Mechelen \work{Bullarium} volume that \emph{is} available
and whose index lists an incipit \term{Non est equidem}
(\url{https://archive.org/details/benedictipapaex00xivgoog}), the
surrounding constitutions are of 1756 and a full-text search of that
volume for \texttt{Guadalup}, \texttt{Mexic} and \texttt{Tepe}
returns nothing. Whether that incipit belongs to the Guadalupan act
of 1754, to a different act of 1756, or to both, could not be settled
here. (3)~Accordingly nothing is asserted in this study about the
brief's wording or recitals, and the frequently repeated association
of the Psalm~147 verse \term{Non fecit taliter omni nationi} with this
act is likewise unverified and not asserted. What \emph{is} verified
is the substance of the operative text quoted above: an oath of 1737
and a papal confirmation of 1754.}

\subsection{The nineteenth century: coronation, revised Office, and
the Council of 1896}

The nineteenth century produced the tradition's greatest public
honours and its gravest crisis of scholarly confidence at the same
time. The two are causally linked, and the sequence is documented in
detail.

In 1884 the maestrescuela of Guadalajara proposed petitioning Rome
for a new Office and Mass for 12 December. The stated motive was
candid: the older Office contained only a hypothetical declaration of
the historical reality of the tradition, which prevented its defenders
from invoking apostolic authority for the events at precisely the
moment when scholars were putting them in
doubt.\endnote{Edmundo O'Gorman, \work{Destierro de sombras},
Appendix~VII, no.~22, read 2026-07-25 at
\url{https://historicas.unam.mx/publicaciones/publicadigital/libros/222c/222c_04_05_07_ApendiceSeptimo.pdf},
citing the \work{\'Album de la coronaci\'on}, I, p.~84: the idea came
from Rafael S.~Camacho, and the purpose was ``obtener la enmienda del
antiguo oficio que s\'olo conten\'ia una declaraci\'on hipot\'etica de
la realidad hist\'orica de la tradici\'on de las apariciones.''}

That the older Roman register really was hypothetical can be checked
independently, because the Sacred Congregation of Rites used it in a
published rescript two years before the revision.

\begin{operativetext}{Sacred Congregation of Rites, rescript
\term{Mexicana}, ``Dubium quoad collocationem statuarum in Ecclesia,''
15 July 1892, extract (Latin, as printed in \work{Acta Sanctae
Sedis})}
\ldots\ ubi B.~M.~V.\ de Guadalupe imago depicta collocanda erit,
medium inter binas marmoreas statuas, scilicet Rmi D\~ni Zumarraga
primi Antistitis Mexicani et cuiusdam Indi Ioannis Didaci, cui Deipara
Virgo \term{fertur apparuisse}.\par
[Responsum:] Simulacra de quibus in casu, collocari poterunt in
Ecclesia, dummodo non exponantur super altare.\endnote{\work{Acta
Sanctae Sedis} 25 (1892--93), pp.~116--117, read 2026-07-25 in the
Holy See's OCR scan
(\url{https://www.vatican.va/archive/ass/documents/ASS-25-1892-93-ocr.pdf}),
pages located and read at first hand; signed Cai.\ Card.\
Aloisi-Masella, Praef., and Ioann.\ Ponzi, Substit. Editorially:
``\ldots where the painted image of the Blessed Virgin Mary of
Guadalupe is to be placed, between two marble statues, namely of the
Most Reverend Lord Zum\'arraga, first bishop of Mexico, and of a
certain Indian, Juan Diego, to whom the Virgin Mother of God
\emph{is said} to have appeared.\ldots\ [Reply:] The statues in
question may be placed in the church, provided they are not exposed
upon the altar.'' The document is a liturgical rescript about statues;
the phrase \term{fertur apparuisse} is incidental to its purpose,
which is exactly what makes it good evidence of the Congregation's
ordinary register. Quoted within ordinary scholarly limits with
citation.}
\end{operativetext}

On 24 September 1886 the three Mexican archbishops petitioned Leo~XIII
to approve the coronation of the image; his brief conceding it is
dated 8 February 1887 and was published with the archbishops' pastoral
letter of 19 March 1887. The petition for the new Office went to Rome
on 15 March 1890 and met sustained resistance: the Congregation of
Rites took it up in 1892 but its prefect raised difficulties, the
promoter of the faith raised further ones, on 15 April 1893 the matter
was deferred \term{cum adnotationibus S.~Fidei Promotoris}, and the
promoter's printed objections reached Mexico at the end of October
1893. An agent's letter of 7 March 1893 reports the prefect saying
that account had to be taken of certain anonymous writings that had
reached the Congregation---almost certainly the clandestine Latin
translation of Garc\'ia Icazbalceta's letter which an opponent of the
tradition had had printed in 1888 and sent to the members of the
Congregation for exactly that purpose.\endnote{O'Gorman, Appendix~VII,
nos.~23, 25--26, 36, 40, 50--54, same witness. The Latin translation
is the anonymous \work{De B.M.V.\ Apparitione in Mexico sub titulo de
Guadalupe, exquisitio historica} (Mexico, 1888), made from a copy of
Garc\'ia Icazbalceta's letter surreptitiously taken from Francisco del
Paso y Troncoso's desk and published by the canon Vicente de
P.~Andrade, who---O'Gorman states---sent copies to the members of the
Congregation of Rites precisely in order to prevent the concession.
Reported at O'Gorman's level; the Roman acta of the cause were not
consulted, and see the recorded negative below.}

The concession was made on 6 March 1894, and the change made in it is
the most exact index available of what Rome was prepared to say. The
word \term{fertur} of the old Office was replaced by the phrase
\term{antiqua et constans traditio docet}, and the Pope approved.

\term{Fertur}---``it is reported''---is the classic liturgical formula
of hagiographical reserve, and the 1892 rescript shows it in live use.
\term{Antiqua et constans traditio docet}---``an ancient and constant
tradition teaches''---is a stronger warrant, and it is still a
statement about a \emph{tradition}. Rome moved one step and stopped
precisely where the historians' argument began.\endnote{O'Gorman,
Appendix~VII, no.~57, same witness: ``Congregaci\'on de cardenales
aprob\'o el nuevo oficio y misa propia en honor de la Virgen de
Guadalupe. Se sustituy\'o la palabra \term{fertur} del antiguo oficio
por la frase \term{antiqua et constante traditio ducet}. El papa
concedi\'o su aprobaci\'on,'' citing \work{Sacra Rituum
Congregatione\ldots} (Rome: Typis Perseverantiae, 1894). O'Gorman's
text prints \term{constante} and \term{ducet}; the standard Latin of
the formula is \term{antiqua et constans traditio docet}, and the
variants are recorded rather than silently corrected. \emph{Recorded
negative search.} The Holy See's own organ of the period was checked
directly on 2026-07-25: \work{Acta Sanctae Sedis} volumes 19
(1886--87), 20 (1887), 26 (1893--94) and 27 (1894--95) were retrieved
and searched for \texttt{Guadalup}; neither the brief of 8 February
1887, nor the decree of 6 March 1894, nor any notice of the 1895
coronation was found in them, and the only Mexican Guadalupan matter
located anywhere in the \work{ASS} of these years is the 1892 rescript
quoted above. The 1894 concession is therefore reported at O'Gorman's
level and at the level of the 1894 Roman printing he cites, which was
not consulted.}

The image was crowned in the collegiate church at ten in the morning
on 12 October 1895.\endnote{O'Gorman, Appendix~VII, nos.~62 and 65,
same witness; the image had been moved to the neighbouring church of
the Capuchins on 24 February 1888 during the works and was returned on
30 September 1895. The date is independently attested fifty years
later by Pius~XII's radio message for the fiftieth anniversary
(\S8.4).} The following year the Fifth Mexican Provincial Council met,
and its prelates issued an edict in response to the recent
publications against the tradition. Its first article is the most
precise statement of the tradition's status that this study located
from any competent ecclesiastical body.

\begin{operativetext}{Prelates of the Fifth Mexican Provincial
Council, edict of 1896, article~1, as reported}
que la maravillosa aparici\'on sin ser un dogma de fe, es una
tradici\'on antigua, constante y universal en la naci\'on
mexicana.\endnote{O'Gorman, Appendix~VII, no.~70, same witness, which
summarizes the edict's six articles and quotes article~1 in these
words. Editorially: ``that the marvellous apparition, without being a
dogma of faith, is an ancient, constant and universal tradition in the
Mexican nation.'' The remaining articles declare that to doubt the
tradition is rash and a grave injury to the prelates who transmitted
it, invoke the pontifical declarations of Benedict~XIV and Leo~XIII,
and order a solemn anniversary of the coronation in reparation. The
Council met 23 August--1 November 1896 and its decrees were reviewed
and approved by the Holy See on 19 August 1898 (\emph{ibid.},
no.~69). \emph{The conciliar acts were not consulted}; this is
reported at O'Gorman's level. The edict is a provincial-council act of
1896, not a determination under any modern procedure, and it is not
relabelled here under the 2024 taxonomy.}
\end{operativetext}

Two things about that formula deserve notice. It is exact: it denies
that the apparition is a dogma and affirms that it is a tradition,
which is what the evidence permits a bishop to say. And it was issued
in a defensive posture, in reparation for publications, which is part
of its meaning and should not be edited out.

\subsection{The twentieth and twenty-first centuries}

\subsubsection*{Patronage of Latin America, and a recorded gap}

On 16 August 1910 the Ordinary Congregation of Sacred Rites met in
the Apostolic Vatican Palace with, as the fifth item of its agenda,
``the election and declaration of the B.~Virgin Mary called of
Guadalupe as Patroness of all Latin America.'' The resulting act was
not located.\endnote{\work{Acta Apostolicae Sedis} 2 (1910), p.~651,
\term{Diarium Romanae Curiae}, read 2026-07-25 in the Holy See's OCR
scan and located at first hand: ``V. Intorno all'elezione e
dichiarazione della B.~Maria Vergine detta di Guadalupe, a Patrona di
tutta l'America Latina.'' \emph{Recorded negative search.} The
complete \work{AAS} volumes 2 (1910) and 3 (1911) were searched on
2026-07-25 for \texttt{uadalup}---a substring chosen so that the
result does not depend on how the initial capital was resolved by the
scan's OCR; \work{AAS} 1910 returns four hits
(Pius~X's letter of 23 February 1910 to the Mexican bishops at
pp.~98--99; his letter of 16 June 1910 at p.~569; the agenda item
above; and an unrelated 1817 reference), and \work{AAS} 1911 returns
none. No decree or brief declaring Our Lady of Guadalupe patroness of
Latin America was published in the \work{Acta Apostolicae Sedis} in
either year, and the commonly cited date of 24 August 1910 is not
verified here.}

\subsubsection*{Pius XII, 1945}

For the fiftieth anniversary of the coronation Pius~XII addressed a
radio message to the Mexican faithful on 12 October 1945. Its register
is worth attending to: the Pope narrates the tradition, and marks
that he is narrating it.

\begin{operativetext}{Pius~XII, radio message to the faithful of
Mexico on the fiftieth anniversary of the canonical coronation of the
Virgin of Guadalupe, 12 October 1945, extracts (Spanish)}
En la tilma del pobrecito Juan Diego ---como refiere la
tradici\'on--- pinceles que no eran de ac\'a abajo dejaban pintada una
imagen dulc\'isima\ldots\par
\guillemotleft\textexclamdown Viva la Virgen de Guadalupe,
Emperatriz de Am\'erica y Reina de M\'ejico!\guillemotright\endnote{Read 2026-07-25 at
\url{https://www.vatican.va/content/pius-xii/es/speeches/1945/documents/hf_p-xii_spe_19451012_guadalupe-mexico.html}.
Editorially: ``On the tilma of poor little Juan Diego---as tradition
relates---brushes that were not from here below left a most sweet
image painted\ldots'' and ``Long live the Virgin of Guadalupe,
Empress of America and Queen of Mexico!'' The guillemets and the
inverted exclamation mark are the Vatican page's own. Note that in the message the
acclamation is first placed in the mouth of the crowd of 1895 and only
afterwards taken up by the Pope in his own voice. \emph{Recorded
negative search.} The complete \work{AAS} volume 38 (1946) was
searched on 2026-07-25 for \texttt{Guadalup}: the only hits concern the
diocese of Guadalupe in Bolivia. No act of Pius~XII declaring Our Lady
of Guadalupe patroness or empress of the Americas was located, and
none is asserted here. Quoted within ordinary scholarly limits with
citation.}
\end{operativetext}

\subsubsection*{The Basilica, 1976}

Two acts of Paul~VI belong to the modern shrine. By the apostolic
letter \term{Sacra illa} of 4 October 1976 the title and dignity of
minor basilica were transferred in perpetuity from the old church to
the newly built one; on 12 October 1976 the new temple was
dedicated.\endnote{\work{Acta Apostolicae Sedis} 68 (1976), p.~598
(\term{Sacra illa}: ``Templum Beatae Mariae Virginis vulgo `de
Guadalupe', Mexicopoli prope antiquum dicatum, titulo ac dignitate
Basilicae minoris exornatur''; granted at the petition of Cardinal
Miguel Dar\'io Miranda y G\'omez, archbishop of Mexico, and dated
\term{die IV mensis Octobris, anno MCMLXXVI}) and pp.~653--655
(Paul~VI's message for the dedication of the new temple, with the
editorial footnote ``Die 12 m.\ octobris a.\ 1976''); read 2026-07-25
in the Holy See's OCR scan
(\url{https://www.vatican.va/archive/aas/documents/AAS-68-1976-ocr.pdf}),
pages located and read at first hand.}

\subsubsection*{John Paul II: pilgrimage, teaching, feast,
canonization}

John Paul~II came to the shrine in 1979, 1990, 1999 and 2002; the last
three visits carried the beatification, the promulgation of
\work{Ecclesia in America}, and the canonization. The exhortation was
given at Mexico City on 22 January 1999, and its paragraph~11 is the
fullest magisterial statement on Guadalupe.

\begin{operativetext}{John Paul~II, \work{Ecclesia in America} (22
January 1999), n.~11 (English), extracts}
The appearance of Mary to the native Juan Diego on the hill of
Tepeyac in 1531 had a decisive effect on evangelization. Its
influence greatly overflows the boundaries of Mexico, spreading to
the whole Continent.\par
\ldots America\ldots\ has recognized in the mestiza face of the Virgin
of Tepeyac, ``in Blessed Mary of Guadalupe, an impressive example of a
perfectly inculturated evangelization''. Consequently, not only in
Central and South America, but in North America as well, the Virgin
of Guadalupe is venerated as Queen of all America.\par
\ldots In the prayer composed for the Special Assembly for America of
the Synod of Bishops, Holy Mary of Guadalupe is invoked as
``Patroness of all America and Star of the first and new
evangelization''. In view of this, I welcome with joy the proposal of
the Synod Fathers that the feast of Our Lady of Guadalupe, Mother and
Evangelizer of America, be celebrated throughout the continent on
December 12.\endnote{English text read 2026-07-25 at
\url{https://www.vatican.va/content/john-paul-ii/en/apost_exhortations/documents/hf_jp-ii_exh_22011999_ecclesia-in-america.html};
the signature clause was located on the same page: ``Given at Mexico
City, January 22, in the year 1999, the twenty-first of my
Pontificate.'' The internal quotations are the exhortation's own,
from the Fourth General Conference of the Latin American Bishops at
Santo Domingo and from the Synod's prayer. Ellipses mark omissions.
The 1999 decree of the Congregation for Divine Worship (below)
distinguishes the promulgation \term{die 22 ianuarii 1999 in Civitate
Mexicana} from the homily \term{postridie in Basilica Dominae Nostrae
de Guadalupe}; this study accordingly does not describe the
exhortation as signed at the Basilica. Note the register of n.~11: it
speaks of \emph{the appearance}, of \emph{recognition} by the peoples
of America, and of a \emph{proposal of the Synod Fathers} which the
Pope welcomes. Its operative content is liturgical.}
\end{operativetext}

The implementing act followed two months later, and it is exact.

\begin{operativetext}{Congregation for Divine Worship and the
Discipline of the Sacraments, Prot.~803/99/L, decree \work{de festo
Beatae Mariae Virginis de Guadalupe in tota America quotannis die 12
decembris peragendo}, 25 March 1999, extracts (Latin)}
\ldots\ per Adhortationem Apostolicam postsynodalem \work{Ecclesia in
America}, die 22 ianuarii 1999 in Civitate Mexicana promulgatam,
atque homiliam postridie in Basilica Dominae Nostrae \term{de
Guadalupe} habitam benigne censuit, ut in universa America celebratio
Beatae Mariae Virginis \term{de Guadalupe} in posterum gradu
\term{festi} peragereture [sic]\ldots\par
Istud igitur Dicasterium hoc festum declarat in calendariis cuiusque
nationis vel territorii Americae ad diem 12 decembris inscribendum,
ut in singulis Continentis dioecesibus eo die quotannis celebretur,
salvis semper concessionibus gradus sollemnitatis ab Apostolica Sede
in favorem quarundam circumscriptionum et ecclesiarum decreto vel
norma iam peractis.\endnote{\work{Notitiae} 35 (1999), issue 394--395,
pp.~272--274; the protocol number, title, text, date (``Ex aedibus
Congregationis de Cultu Divino et Disciplina Sacramentorum, in
Sollemnitate Annuntiationis Domini, die 25 martii 1999'') and
signatures (Georgius A.\ Card.\ Medina Est\'evez, \term{Praefectus};
Marius Marini, \term{Subsecretarius}) were read at first hand on
2026-07-25 from rendered page images of the Dicastery's own scan at
\url{https://www.cultodivino.va/content/dam/cultodivino/rivista-notitiae/1990/notitiae-35-(1999)/Notitiae-394-395-1999.pdf}
(printed pages 272, 273 and 274). The bracketed \emph{sic} marks the
reading of the printed page. Editorially: ``\ldots he graciously
determined, through the post-synodal apostolic exhortation
\work{Ecclesia in America}, promulgated on 22 January 1999 in Mexico
City, and through the homily delivered the following day in the
Basilica of Our Lady of Guadalupe, that in the whole of America the
celebration of the Blessed Virgin Mary of Guadalupe should henceforth
be kept with the rank of a feast\ldots\ This Dicastery therefore
declares that this feast is to be inscribed on 12 December in the
calendars of each nation or territory of America, so that it may be
celebrated on that day each year in every diocese of the Continent,
always saving the concessions of the rank of solemnity already granted
by the Apostolic See by decree or norm in favour of certain
circumscriptions and churches.'' The decree also calls her
\term{totius Americae Imperatrix} and declares the annexed Spanish
texts \term{typici}. Quoted within ordinary scholarly limits with
citation.}
\end{operativetext}

On 31 July 2002, in the Basilica, John Paul~II canonized Juan Diego
Cuauhtlatoatzin (\S5). The homily's register is again worth notice: it
calls the Guadalupe writings ``the accounts of Guadalupe, sensitively
written and steeped in tenderness''; it quotes the Mexican episcopate's
statement of 14 May 2002 that the ``Guadalupe Event'' meant a
beginning of evangelization ``with a vitality that surpassed all
expectations''; it describes Juan Diego as having ``facilitated the
fruitful meeting of two worlds''; and it turns at once to a
present-tense obligation: ``it is necessary today to support the
indigenous peoples in their legitimate aspirations, respecting and
defending the authentic values of each ethnic
group.''\endnote{John Paul~II, homily at the canonization of Juan
Diego Cuauhtlatoatzin, Mexico City, 31 July 2002, nn.~3--4, English
text read 2026-07-25 at
\url{https://www.vatican.va/content/john-paul-ii/en/homilies/2002/documents/hf_jp-ii_hom_20020731_canonization-mexico.html};
the Spanish text is in \work{AAS} 94 (2002), p.~743ff. The homily's
closing words to the new saint---``Show us the way that leads to the
`Dark Virgin' of Tepeyac, that she may receive us in the depths of
her heart, for she is the loving, compassionate Mother who guides us
to the true God''---are a compact statement of the Marian boundary of
\S9.1. Quoted within ordinary scholarly limits with citation.}

\subsubsection*{The General Roman Calendar, 2002 and 2008}

One liturgical fact is regularly misstated and is worth fixing. The
third typical edition of the \work{Missale Romanum} (2002) did not
carry Our Lady of Guadalupe or St Juan Diego in the General Roman
Calendar. Both were added in the emended reprint of 2008, whose
official list of variations records, under
\term{Calendarium Romanum Generale}, the addition of
\term{S.~Ioannis Didaci Cuauhtlatoatzin} at 9 December and of
\term{B.~Mariae Virginis de Guadalupe} at 12
December.\endnote{Congregation for Divine Worship and the Discipline
of the Sacraments, \work{Variationes et additiones in reimpressione
emendata ``Missalis Romani,'' editionis typicae tertiae},
\work{Notitiae} 44 (2008), issue 503--504, under
\term{Calendarium Romanum Generale} (``p.~116: Ad diem 9: additur:
`S.~Ioannis Didaci Cuauhtlatoatzin'.'' ``p.~116: Ad diem 12: additur:
`B.~Mariae Virginis de Guadalupe'.''). The absence of both from the
2002 edition is established from the Dicastery's own presentation of
the third typical edition, \work{Notitiae} 38 (2002), pp.~457--459,
which enumerates the celebrations added to the General Roman Calendar
and includes neither. \emph{Evidentiary ceiling}: these two
\work{Notitiae} loci are reported at the level of the verification
performed for this study by an assisting research agent working from
the Dicastery's own scans at \url{https://www.cultodivino.va/}; unlike
the 1999 decree, they were not re-read page by page by the author of
this study, and they are marked accordingly. \emph{Recorded limit}:
the 2008 list adds the two entries without a rank qualifier. This
study does not assert the rank of either celebration in the General
Roman Calendar, and does not assert the rank in any particular
calendar; the 1999 decree's \term{gradu festi} governs the calendars
of the Americas, and the United States Conference of Catholic Bishops
presents 12 December as a Feast.}

\subsubsection*{Continuing reception}

Papal celebration of the feast at Rome has continued: Francis
celebrated Masses of Our Lady of Guadalupe in St Peter's Basilica on
12 December in several years and went on pilgrimage to the Basilica in
Mexico City on 13 February 2016; Leo~XIV preached at the Mass of
Our Lady of Guadalupe in St Peter's on 12 December 2025 and, on 5
February 2026, addressed a message to a pastoral-theological congress
on the Guadalupan event, looking toward the fifth centenary of the
apparition in 2031.\endnote{Verified 2026-07-25 by an assisting
research agent against vatican.va document trees, and reported at that
level: Francis, homily at the Basilica of Our Lady of Guadalupe,
Mexico City, 13 February 2016
(\url{https://www.vatican.va/content/francesco/en/homilies/2016/documents/papa-francesco_20160213_omelia-messico-guadalupe.html});
homilies for 12 December present on vatican.va for 2019, 2020, 2022,
2023 and 2024 (no Guadalupe homily is indexed for 2021); Leo~XIV,
homily, ``Blessed Virgin Mary of Guadalupe --- Holy Mass,'' St Peter's
Basilica, 12 December 2025
(\url{https://www.vatican.va/content/leo-xiv/en/homilies/2025/documents/20251212-messa-guadalupe.html});
Leo~XIV, message to participants in the pastoral theological congress
on the Guadalupan event, 5 February 2026
(\url{https://www.vatican.va/content/leo-xiv/en/messages/pont-messages/2026/documents/20260205-messaggio-congresso-guadalupe.html}).
These are acts of papal piety and teaching in their proper genres;
none is a determination of the historical question.}

\subsection{What the chain of acts amounts to}

Set out in order, the acts form a consistent picture, and it is not
the picture that either party in the historical dispute would most
like.

The Church has, over three centuries, confirmed a national patronage,
conceded a proper Office and Mass, revised its wording once under
scrutiny and in a direction that went one step and no further, crowned
the image, transferred basilica dignity to a new shrine built for the
crowds, extended a continental feast of the rank of \term{festum},
inscribed the celebration in the General Roman Calendar, canonized the
reported seer, and sent four papal pilgrimages to Tepeyac. That is an
extraordinary weight of reception, and any account that treats
Guadalupe as merely tolerated is false to the record.

At no point in that chain is there an act determining that the events
of 1531 occurred. The nearest approaches are a liturgical formula
asserting an ancient and constant \emph{tradition}; a provincial
council's edict expressly denying that the apparition is a dogma of
faith; and an apostolic letter of beatification which narrates the
event as ``the writings relate'' and states in its own text that not
many notices about Juan Diego have come down to us. Any account that
treats Guadalupe as juridically established history is equally false
to the record.

Both statements are true at once, and holding them together is the
whole discipline of this study.
